\par\bigskip
\SourceLeaf{424}
\begin{center}
\pdfbookmark[2]{Gauss an Olbers, 13. December 1810}{pallas-brief-18101213}
\textsc{Gauss} an \textsc{Olbers}, 13. December 1810.
\end{center}
\noindent\hspace*{\parindent}Meine beiden letzten Briefe haben sich mit dem Ihrigen gekreuzt. Ich habe seitdem meine Arbeit über die Störungen der Pallas durch Jupiter vollendet, und eine meine kühnste Erwartung übertreffende Übereinstimmung herausgebracht. Sehen Sie hier die Unterschiede:
\begin{center}\begingroup\small\setlength{\tabcolsep}{4pt}
\begin{tabular}{lrrrrrr}
& 1803 & 1804 & 1805 & 1807 & 1808 & 1809\\
Mittl. Länge & \(+1''{,}3\) & \(-3''{,}8\) & \(+3''{,}9\) & \(-3''{,}3\) & \(+3''{,}2\) & \(-1''{,}4\)\\
Helioc. Breite & \(-1''{,}0\) & \(+4''{,}1\) & \(+6''{,}0\) & \(+3''{,}9\) & \(-16''\) & \(-3''{,}9.\)\\
\end{tabular}\endgroup\end{center}
\par
\noindent\hspace*{\parindent}Bei Berechnung der Störungen hatte ich die Elemente zum Grunde gelegt, die Sie auf den beiliegenden Blättern\HistoricalNoteEight{Band VI, S. 63.]} pag. [1971] als No. IV finden, ich habe angefangen, die ganze Rechnung nochmals zu wiederholen, indem ich den ganzen Zeitraum von 1803—1811 in sieben oder acht Perioden theile und in jeder besondere Elemente zu Grunde lege, wie sie meine erste Rechnung gegeben hat. Aber wie werden Sie mit mir erstaunen, wenn ich Ihnen sage, dass meine Ephemeride im Oct. der M. C., wenn ich mich auf die Resultate der Störungsrechnungen verlassen darf, jetzt über \(1^{0}\) falsch ist? Ich glaube zwar meiner Sache gewiss zu sein, doch bin ich, da die Rechnung so sehr weitläuftig ist, nicht eher ganz ruhig, als bis Pallas wiedergefunden ist, wozu ich jetzt mit Ungeduld günstiges Wetter wünsche. Sollten Sie in Bremen glücklicher sein, so sehen Sie sich auch wohl einmal danach um. Nach den gestörten Elementen sollte die AR den 19. Dec. um \(1^{0}10'\), den 10. Jan. um \(1^{0}40'\) kleiner, die südliche Declination um \(12'\,.\,.\,.\,10'\) grösser sein als in der Ephemeride. In der {\AstroSymbols ☌} 1809 mit {\AstroSymbols ♃} ist {\AstroSymbols ⚴} gar zu sehr gestört. Bestätigt sich meine Rechnung auch hier zu meiner Zufriedenheit, so werde ich vielleicht Juno und Vesta auf eine ähnliche Art behandeln.
\par
\noindent\hspace*{\parindent}Die Abhandlung, wovon ich Ihnen hier die Anzeige beilege, wird bald gedruckt werden.\quad.\quad.\quad.\quad.\quad.
\par
